\chapter{متغیرها و ریخت‌های داده}%
\label{ch:variables}

در فصلِ پیش با متغیر آشنا شدیم و آن را به جعبه‌ای تشبیه کردیم که چیزی در آن
نگه می‌داریم. در این فصل ژرف‌تر می‌شویم: ریخت‌های گوناگونِ داده را می‌شناسیم،
یاد می‌گیریم متغیر را چگونه بسازیم، و درمی‌یابیم چرا «ریختِ» هر متغیر اهمیت
دارد.

\section{متغیر؛ جعبه‌ای با نام}

تصور کنید چند جعبه دارید و روی هر کدام برچسبی زده‌اید: «سن»، «قد»، «نام». در هر
جعبه چیزی گذاشته‌اید. هر وقت به محتوای جعبه نیاز داشتید، کافی است نامِ روی برچسب
را صدا بزنید. در برنامه‌نویسی، این جعبه‌ها همان \textbf{متغیر}ها هستند.

ساده‌ترین شکلِ ساختنِ یک متغیر در سلام چنین است:

\begin{salamfa}
کارکرد آغازین:
    گذرا سن : صحیح = 30
    چاپ سن
پایان
\end{salamfa}

این خط را چنین بخوانید: «یک \textbf{متغیر} به نامِ \textbf{سن} که ریختش
\textbf{صحیح} (عددِ درست) است بساز و مقدارِ \textbf{۳۰} را در آن بگذار.»

ساختار همیشه یکسان است:

\begin{center}
\code{گذرا} \;\textit{نام}\;:\;\textit{ریخت} \;=\; \textit{مقدار}
\end{center}

\section{چهار شیوهٔ ساختنِ متغیر}

سلام چهار واژهٔ کلیدی برای ساختنِ متغیر دارد و هر کدام معنایی متفاوت دارند. به
این برنامهٔ کوچک که تقریباً همان نمونهٔ موجود در پوشهٔ مثال‌های سلام است ---
نگاه کنید:

\begin{salamfa}
کارکرد آغازین:
    گذرا سن : صحیح = 30
    پایدار سال تولد : صحیح = 1994
    قد : اعشار۶۴ = 1.75
    نام : خودکار = "سارا"
    سن = سن + 1
    چاپ نام, سن, سال تولد, قد
پایان
\end{salamfa}

بیایید چهار شیوه را از هم تفکیک کنیم:

\subsection{متغیر: جعبه‌ای که می‌توان عوضش کرد}

\begin{salamfa}
گذرا سن : صحیح = 30
\end{salamfa}

واژهٔ \code{گذرا} (معادلِ \code{mut} در انگلیسی) جعبه‌ای می‌سازد که محتوایش را
بعداً می‌توان \emph{تغییر} داد. به همین دلیل در خطِ \code{سن = سن + 1}
توانستیم مقدارِ \code{سن} را یکی زیاد کنیم؛ حالا \code{سن} برابرِ ۳۱ است.

\subsection{پایدار: جعبه‌ای که قفل است}

\begin{salamfa}
پایدار سال تولد : صحیح = 1994
\end{salamfa}

واژهٔ \code{پایدار} (معادلِ \code{const}) جعبه‌ای می‌سازد که پس از پُر شدن،
\emph{دیگر تغییر نمی‌کند}. سالِ تولدِ کسی هیچ‌گاه عوض نمی‌شود، پس درست است که
آن را پایدار بگیریم. اگر بکوشید مقدارِ یک پایدار را تغییر دهید، کامپایلر خطا
می‌گیرد و به شما هشدار می‌دهد.

\subsection{اعلانِ ساده: پایدار اما بی‌واژه}

\begin{salamen}
قد : اعشار۶۴ = 1.75
\end{salamen}

اگر هیچ واژه‌ای (\code{گذرا} یا \code{پایدار}) ننویسید و مستقیم با نام آغاز
کنید، یک متغیرِ \emph{تغییرناپذیر} ساخته می‌شود درست مانندِ \code{پایدار}.
این شکلِ کوتاه و تمیز، شیوهٔ پیش‌فرضِ سلام است: \emph{مگر آنکه واقعاً نیاز به
تغییر داشته باشید، متغیر را تغییرناپذیر بگذارید.}

\subsection{خودکار: بگذارید سلام ریخت را حدس بزند}

\begin{salamen}
نام : خودکار = "سارا"
\end{salamen}

گاهی نوشتنِ ریخت، تکراری و آزاردهنده است؛ مگر معلوم نیست که \code{"سارا"} یک
متن است؟ واژهٔ \code{خودکار} (معادلِ \code{auto}) به سلام می‌گوید: «خودت از روی
مقدار، ریخت را تشخیص بده.» چون \code{"سارا"} میانِ نقل‌قول است، سلام درمی‌یابد
که ریختش \code{رشته} است.

\begin{tip}
استفاده از \code{خودکار} هم کد را کوتاه‌تر می‌کند و هم خواناتر؛ اما در جاهایی
که ریخت روشن نیست (مثلاً یک عددِ خام که نمی‌دانید درست است یا اعشاری)، نوشتنِ
صریحِ ریخت به خوانندهٔ کدِ شما کمک می‌کند.
\end{tip}

\begin{note}
به شیوهٔ صدا زدنِ \code{چاپ} دقت کنید: آرگومان‌ها را \textbf{بدونِ پرانتز} و
جدا شده با ویرگول می‌نویسیم، مانندِ \code{چاپ "سلام،", نام} یا
\code{چاپ نام, سن}. اگر همهٔ آرگومان‌ها را در پرانتز بپیچید، سلام خطا
می‌گیرد؛ \code{چاپ} یک دستور است، نه یک کارکردِ معمولی.
البته پرانتز در دلِ خودِ آرگومان‌ها آزاد است و برای گروه‌بندیِ محاسبات به کار
می‌رود؛ مثلاً \code{چاپ 1, (3 * 4)} یا \code{چاپ (2 + 3) * 4} کاملاً درست‌اند.
\end{note}

\section{ریخت‌های پایهٔ داده}

هر متغیر یک \textbf{ریخت} دارد که می‌گوید چه جور داده‌ای را می‌تواند نگه دارد.
ریخت، هم به سلام کمک می‌کند که برنامهٔ شما را درست بفهمد، و هم از بسیاری خطاها
جلوگیری می‌کند. مهم‌ترین ریخت‌های پایه در سلام این‌ها هستند:

\subsection{اعدادِ درست (صحیح)}

عددهای بدونِ اعشار مانندِ ۰، ۷، ۱۵، و ۱۰۰- را اعدادِ \emph{درست} یا
\emph{صحیح} می‌نامیم. سلام چند ریختِ عددِ درست دارد که در \emph{اندازه} با هم
فرق می‌کنند:

\begin{center}
\begin{tabular}{c c l}
\toprule
\textbf{نامِ فارسی} & \textbf{نامِ انگلیسی} & \textbf{توضیح} \\
\midrule
\code{صحیح} & \code{int} / \code{i32} & عددِ درستِ معمول (۳۲ بیتی) \\
\code{صحیح۳۲} & \code{i32} & عددِ درستِ ۳۲ بیتی \\
\code{صحیح۶۴} & \code{i64} & عددِ درستِ ۶۴ بیتی (برای عددهای خیلی بزرگ) \\
\bottomrule
\end{tabular}
\end{center}

برای کارهای روزمره \code{صحیح} کافی است. اگر با عددهای بسیار بزرگ سر و کار
دارید (مثلاً فاکتوریلِ اعداد یا شمارشِ چیزهای فراوان)، از \code{صحیح۶۴}
استفاده کنید.

\begin{note}
\code{صحیح} و \code{صحیح۳۲} دقیقاً \emph{یک} ریخت‌اند: هر دو همان عددِ درستِ ۳۲
بیتی (\code{i32}) هستند و هر جا یکی را بنویسید، دیگری هم همان معنا را می‌دهد.
عددِ درست در سلام به‌طورِ پیش‌فرض ۳۲ بیتی است، پس نوشتنِ \code{۳۲} در پایانِ
نام ضروری نیست. در سراسرِ این کتاب همان شکلِ کوتاه، یعنی \code{صحیح}، را به‌کار
می‌بریم؛ کوتاه‌تر است و کد را خواناتر نگه می‌دارد. \code{صحیح۳۲} را تنها وقتی
بنویسید که بخواهید بر پهنای ۳۲ بیتیِ عدد تأکید کنید.
\end{note}

\begin{salamfa}
کارکرد آغازین:
    گذرا کوچک : صحیح = 1000
    گذرا بزرگ : صحیح۶۴ = 9000000000
    چاپ کوچک, بزرگ
پایان
\end{salamfa}

\subsection{اعدادِ اعشاری}

عددهایی که بخشِ اعشاری دارند مانندِ ۱٫۷۵، ۳٫۱۴ و ۰٫۵ را اعدادِ
\emph{اعشاری} می‌نامیم. در سلام ریختِ آن‌ها \code{اعشار۶۴} است (معادلِ
\code{f64}):

\begin{salamfa}
کارکرد آغازین:
    قد : اعشار۶۴ = 1.75
    وزن : اعشار۶۴ = 68.5
    چاپ قد, وزن
پایان
\end{salamfa}

\begin{warn}
در عددهای اعشاری، نقطهٔ اعشار همیشه با نقطهٔ لاتین (\code{.}) نوشته می‌شود، نه
ممیزِ فارسی. پس بنویسید \code{1.75} نه «۱٫۷۵». همچنین برای آنکه عددی اعشاری
شمرده شود، باید نقطه داشته باشد: \code{2} یک عددِ درست است اما \code{2.0} یک
عددِ اعشاری.
\end{warn}

\subsection{متن (رشته)}

هر دنباله‌ای از نویسه‌ها یک واژه، یک جمله، یا حتی یک نام یک
\emph{رشته} (string) است. رشته‌ها را همیشه میانِ دو نقل‌قول می‌نویسیم و ریختشان
\code{رشته} است (معادلِ \code{str}):

\begin{salamfa}
کارکرد آغازین:
    شهر : رشته = "اصفهان"
    پیام : رشته = "به سلام خوش آمدید"
    چاپ شهر, پیام
پایان
\end{salamfa}

\subsection{مقادیرِ منطقی (درست/نادرست)}

گاهی فقط دو پاسخ ممکن است: آری یا نه، روشن یا خاموش، درست یا نادرست. برای این
موارد ریختِ \code{منطقی} (معادلِ \code{bool}) را داریم که تنها دو مقدار می‌پذیرد:
\code{درست} (true) و \code{نادرست} (false):

\begin{salamfa}
کارکرد آغازین:
    باران : منطقی = درست
    آفتاب : منطقی = نادرست
    چاپ باران, آفتاب
پایان
\end{salamfa}

مقادیرِ منطقی هنگامِ تصمیم‌گیری بسیار به‌کار می‌آیند؛ در فصلِ شرط‌ها (فصلِ~\ref{ch:conditionals})
نقشِ کلیدیِ آن‌ها را خواهیم دید.

\subsection{کاراکتر (تک‌نویسه)}

اگر تنها به \emph{یک} نویسه نیاز داشته باشید مثلاً حرفِ \lr{A} یا رقمِ ۷ ---
از ریختِ \code{کاراکتر} (معادلِ \code{char}) استفاده می‌کنید. کاراکترها را میانِ
\emph{تک‌نقل‌قول} می‌نویسیم:

\begin{salamfa}
کارکرد آغازین:
    حرف : کاراکتر = 'A'
    چاپ حرف
پایان
\end{salamfa}

\begin{note}
تفاوتِ \code{'A'} و \code{"A"} چیست؟ نخستی یک \emph{کاراکتر} (تک‌نویسه) است و
دومی یک \emph{رشته} (دنباله‌ای از نویسه‌ها که اینجا فقط یک نویسه دارد). در
بسیاری از زبان‌ها این تمایز مهم است.
\end{note}

\section{جدولِ کاملِ ریخت‌های پایه}

\begin{center}
\begin{tabular}{c c l}
\toprule
\textbf{فارسی} & \textbf{انگلیسی} & \textbf{چه چیزی نگه می‌دارد} \\
\midrule
\code{صحیح} / \code{صحیح۳۲} & \code{int} / \code{i32} & عددِ درستِ ۳۲ بیتی \\
\code{صحیح۶۴} & \code{i64} & عددِ درستِ ۶۴ بیتی \\
\code{اعشار۶۴} & \code{f64} & عددِ اعشاری \\
\code{اعشار۳۲} & \code{f32} & عددِ اعشاریِ کوچک‌تر \\
\code{منطقی} & \code{bool} & درست یا نادرست \\
\code{کاراکتر} & \code{char} & یک نویسه \\
\code{رشته} & \code{str} & دنباله‌ای از نویسه‌ها \\
\code{خودکار} & \code{auto} & ریخت را خودِ سلام تشخیص می‌دهد \\
\bottomrule
\end{tabular}
\end{center}

\begin{tip}
هر پرونده در سلام به \emph{یک} زبان نوشته می‌شود و نام‌های ریخت هم باید به همان
زبان باشند. در یک برنامهٔ فارسی باید بنویسید \code{گذرا سن: صحیح = 30}؛ اگر
نامِ انگلیسیِ همان ریخت، یعنی \code{i32}، را در پروندهٔ فارسی به‌کار ببرید، کامپایلر
خطا می‌گیرد و می‌گوید نامِ ریخت در یک پروندهٔ فارسی باید فارسی باشد. نامِ انگلیسیِ
ریخت‌ها تنها در پرونده‌های انگلیسی معتبر است.
\end{tip}

\section{تبدیلِ ریخت با «برگردان»}

گاهی می‌خواهید یک مقدار را از ریختی به ریختِ دیگر تبدیل کنید؛ مثلاً یک عددِ درست
را به اعشاری، یا برعکس. برای این کار از واژهٔ \code{برگردان} (معادلِ \code{as})
استفاده می‌کنید:

\begin{salamfa}
کارکرد آغازین:
    گذرا شمار : صحیح = 7
    گذرا دقیق : اعشار۶۴ = شمار برگردان اعشار۶۴
    چاپ دقیق
پایان
\end{salamfa}

اینجا \code{شمار} (که یک عددِ درست است) به \code{اعشار۶۴} تبدیل شد. به این کار
\textbf{تبدیلِ ریخت} (cast) می‌گویند و در فصل‌های بعد به‌ویژه هنگامِ کار با
شمارشی‌ها و آرایه‌ها بارها به آن نیاز خواهیم داشت.

\section{قاعده‌های نام‌گذاری}

نامِ متغیرها را خودتان برمی‌گزینید، اما چند قاعدهٔ ساده هست:

\begin{itemize}
  \item نام می‌تواند فارسی، انگلیسی یا ترکیبی باشد: \code{سن}، \code{age}،
        \code{شمارهٔ‌حساب}.
  \item نام می‌تواند چندواژه‌ای باشد؛ بینِ واژه‌ها فاصله بگذارید:
        \code{سال تولد}، \code{شمارنده اصلی}. (زیرخط \code{\_} هم مجاز است.)
  \item نام نمی‌تواند با یک رقم آغاز شود.
  \item نام نباید یکی از واژگانِ کلیدیِ زبان (مانندِ \code{کارکرد} یا
        \code{اگر}) باشد.
\end{itemize}

\begin{tip}
نامِ خوب، نامی است که \emph{معنا}ی متغیر را برساند. \code{مجموعِ‌نمرات} بسیار
گویاتر از \code{م} است. کدی که شش ماه بعد بخوانید، باید همچنان قابلِ فهم باشد ---
و نامِ خوب نیمی از این فهم را تأمین می‌کند.
\end{tip}

\begin{deep}[نگاهی به درون: ریخت چگونه به ایمنی کمک می‌کند]
سلام یک زبانِ \emph{ایستا-ریخت} (statically typed) است؛ یعنی ریختِ هر متغیر پیش
از اجرا و در زمانِ کامپایل مشخص است. این ویژگی به کامپایلر اجازه می‌دهد بسیاری
از خطاها را \emph{پیش} از اجرای برنامه بگیرد. برای نمونه، اگر بکوشید یک رشته را
با یک عدد جمع بزنید، یا مقدارِ اعشاری را در جعبه‌ای از ریختِ صحیح بگذارید،
کامپایلر همان لحظه به شما هشدار می‌دهد به‌جای آنکه برنامه هنگامِ اجرا و
شاید در دستِ کاربر، ناگهان از کار بیفتد.
\end{deep}

\section{خلاصهٔ فصل}

\begin{itemize}
  \item متغیر، جعبه‌ای نام‌دار برای نگه‌داشتنِ داده است.
  \item چهار شیوهٔ اعلان داریم: \code{گذرا} (تغییرپذیر)، \code{پایدار}
        (تغییرناپذیر)، اعلانِ ساده (تغییرناپذیر)، و \code{خودکار} (تشخیصِ
        خودکارِ ریخت).
  \item ریخت‌های پایه: \code{صحیح}، \code{صحیح۶۴}، \code{اعشار۶۴}،
        \code{منطقی}، \code{کاراکتر} و \code{رشته}.
  \item با \code{برگردان} می‌توان ریختِ یک مقدار را تبدیل کرد.
  \item نام‌های فارسی و انگلیسیِ ریخت‌ها هر دو معتبرند.
\end{itemize}

\section{تمرین‌ها}

\exercise سه متغیر بسازید: نامِ شهر (رشته)، جمعیتِ آن (صحیح۶۴) و میانگینِ دمای
سالانه (اعشار۶۴). سپس هر سه را در یک خط چاپ کنید.

\exercise یک پایدار به نامِ \code{پی} با مقدارِ \code{3.14} بسازید و کوشش کنید
در خطِ بعد مقدارش را تغییر دهید. کامپایلر چه می‌گوید؟

\exercise متغیری از ریختِ \code{منطقی} بسازید که نشان دهد امروز تعطیل است یا
نه، و آن را چاپ کنید.

\exercise عددِ درستِ \code{9} را در یک متغیر بگذارید، سپس آن را \code{برگردان
اعشار۶۴} تبدیل کرده و در متغیرِ دیگری بگذارید و هر دو را چاپ کنید.

\exercise برنامه‌ای بنویسید که از \code{خودکار} برای ساختنِ سه متغیر (یک عدد،
یک اعشاری و یک رشته) استفاده کند و هر سه را چاپ کند.
